\documentclass[12pt]{article}

\usepackage{problem_set_style}
\usepackage{kotex}
\usepackage{booktabs}
\usepackage{array}
\usepackage{graphicx}

\begin{document}
\heading{16}{2027. X. XX (예정)}{}{Spring 2027}{기말 프로젝트 공고 (샘플)}

\section*{1. 프로젝트 개요}
9$\sim$15주차에서 배운 딥러닝 기법(신경망 기초·역전파, CNN,
시퀀스 모델, 어텐션·트랜스포머, 그래프 표현학습, 표현학습, 잠재변수
생성모델)을 실제 데이터에 적용하는 학기 후반 8주에 걸쳐 진행하는 팀 프로젝트이다. 중간 프로젝트
(8주차)의 절차를 \textbf{딥러닝 무대}에서 다시 수행하며, 이번에는
``왜 고전 머신러닝이 아니라 딥러닝인가''를 직접 증명해야 한다.

\vspace{4pt}
\noindent 진행 순서: \textbf{데이터 선정 $\to$ Abstract $\to$ 구현 $\to$
검증 $\to$ 보고서 제출 $\to$ 동료평가}.

\section*{2. 필수 요구사항}
\begin{itemize}
\item \textbf{데이터}: Kaggle, UCI, Hugging Face Datasets 중에서
  선택 (이미지·텍스트·시퀀스 등 비정형 데이터도 가능).
\item \textbf{적용 알고리즘}: 9$\sim$15주차 중 \textbf{최소
  하나의 딥러닝 기법}(CNN, RNN/Transformer, 프롬프팅 응용, 또는
  EM/VAE 계열 생성모델)을 실제 데이터에 적용한다.
\item \textbf{수식적 정당화 최소 1개}: (1) CNN 구조라면 파라미터
  수/연산량 공식으로 설계를 정당화, (2) 정규화(dropout 등)를
  썼다면 그 효과를 설명, (3) 생성모델이라면 ELBO나 min-max
  목적함수를 실제 구현과 연결.
\item \textbf{GBDT 비교 (필수)}: 같은 데이터에 GBDT(또는 랜덤
  포레스트)를 \textbf{베이스라인}으로 돌려, 신경망이 이보다
  나은지 확인한다. 이것이 보고서에서 ``왜 딥러닝인가''를 답하는
  근거의 출발점이 된다.
\item \textbf{검증 원칙}: train/validation/test 3분할을 중간 프로젝트와
  동일하게 준수하고, \textbf{테스트 성능은 딱 한 번}만 확인한다.
\end{itemize}

\newpage
\section*{3. 정형/비정형 판정 가이드}
중간 프로젝트의 5개 데이터 선정 기준(행의 물리적 단위, 레이아웃 함정,
시계열, 클래스 비율, 크기)은 그대로 적용되며, 기말 프로젝트에서는 다음
기준이 추가된다: \textbf{``원시 특징''이 이미 정형으로 추출되어
있는가?}
\begin{center}
\small
\begin{tabular}{@{}p{3.0cm}p{5.2cm}p{5.5cm}@{}}
\toprule
경우 & 예 & 판정 \\
\midrule
수치·범주 열로만 구성된 표 & 매출표, 센서 로그 & \textbf{정형}
--- GBDT와 비교해야 한다 \\
픽셀·음파·토큰인 원시 배열 & 이미지, 오디오, 문장 &
\textbf{비정형} --- 손으로 특징을 만들기 어려우니 신경망이 특징
추출기 역할을 한다 \\
애매한 경우 & 이미 임베딩으로 벡터화된 텍스트 & 원시 토큰으로
직접 모델링하는 선택도 가능 --- 그 이유를 Abstract에 적어야 한다 \\
\bottomrule
\end{tabular}
\end{center}

\vspace{4pt}
\noindent\textbf{데이터 가게와 데이터 월드컵도 활용한다.} \texttt{smhanlab.com/book-ml/datashop}에는 이미지·텍스트·시계열·그래프처럼 딥러닝이 어울리는 현장 데이터(품질 검사 영상, 한국어 리뷰, 음성, 위성 영상, 링크 예측, 차량 궤적 등)도 진열되어 있다. 데이터마다 적힌 행의 단위, 라벨, 크기, 함정을 위 기준으로 먼저 점검하고, 후보를 정하기 어려우면 \textbf{데이터 월드컵}(\texttt{.../datashop/worldcup/})으로 후보를 2$\sim$3개로 좁힌다. 정형으로 쓸 수 있는 데이터를 골랐다면 GBDT 베이스라인과의 비교가 더 중요해진다.

\section*{4. 일정 (후반 8주)}
\begin{center}
\small
\begin{tabular}{@{}p{3.6cm}p{5.4cm}p{3.7cm}@{}}
\toprule
\textbf{시기} & \textbf{해야 할 일} & \textbf{제출} \\
\midrule
\textbf{기말 1주(9주차)} & 팀 구성, 데이터 후보 2$\sim$3개 조사, 정형/비정형 판정 &  \\
\textbf{기말 2주(10주차)} & 행·라벨·지표·베이스라인, ``왜 딥러닝인가'' 초안, 적용할 신경망 계획 & \textbf{Abstract} (1p) \\
\textbf{기말 3$\sim$5주} & 3분할 스캐폴드, train으로만 fit한 전처리, \textbf{GBDT 베이스라인}, 신경망 적용, validation 튜닝, \textbf{test 최종 측정(한 번만)}, 학습 곡선(에포크 단위) &  \\
\textbf{기말 6주(14주차)} & 보고서 완성 & \textbf{보고서 제출}: 보고서 PDF + 코드 + 학습 곡선 + GBDT 비교 \\
\textbf{기말 7주(15주차)} & 제출된 다른 팀 보고서 읽기 &  \\
\textbf{기말 8주(16주차, 기말고사 기간)} & 다른 팀 2개에 대한 peer-review 작성 & \textbf{동료평가} 제출 (리뷰 2건) \\
\bottomrule
\end{tabular}
\end{center}
\vspace{4pt}
학습 곡선은 \textbf{에포크} 단위(중간 프로젝트의 트리 개수와는
다른 축)이다. 최종 숫자만으로는 에포크 수를 어떻게 정했는지
검증할 수 없다 --- validation 곡선이 실제로 쓰였는지가 그림에서
드러나야 한다.

\section*{5. 보고서 구조}
중간 프로젝트의 6절 구조를 그대로 쓰되, 2절과 5절이 기말 프로젝트용으로
바뀐다.
\begin{enumerate}
\item \textbf{문제 정의} (0.5p): 행의 단위, 라벨, 주 지표,
  \textbf{두 층의 베이스라인}(무작위·GBDT), \textbf{정형/비정형
  판정과 그 근거}.
\item \textbf{데이터와 전처리} (1p): 원본 통계, train으로만 fit한
  전처리 명시(이미지 스케일링, 텍스트 토크나이저/패딩 등), 분할
  비율.
\item \textbf{모델 선택과 이유} (1$\sim$2p): 적용한 딥러닝 기법
  1개 이상과 그 구조(층·크기)를 파라미터 수 공식으로 정당화, 왜
  GBDT가 아니라 이 기법인가.
\item \textbf{결과} (1p): 최종 test 숫자(한 번), 에포크 학습 곡선
  1장, GBDT·무작위와의 비교표.
\item \textbf{수식적 정당화} (1p, \textbf{필수 1개}): 파라미터
  수/정규화 효과/ELBO 중 하나.
\item \textbf{잘 안 된 부분과 원인 진단} (0.5p, \textbf{필수}):
  학습이 안 된/과적합된 부분을 배운 개념으로 진단.
\end{enumerate}

\section*{6. 채점 기준 (Rubric)}
\begin{center}
\small
\begin{tabular}{@{}lcp{6.8cm}@{}}
\toprule
항목 & 비중 & 좋은 답 \\
\midrule
문제 정의와 데이터 & 15\% & 행·라벨·지표·\textbf{두 층 베이스라인}이
명확하고, \textbf{정형/비정형 판정}이 데이터 특성에 근거함 \\
방법론 적절성 & 25\% & \textbf{GBDT와 비교해서} 왜 이 딥러닝
기법인지의 근거가 차원·구조·가중치 공유 등의 개념으로 설명됨 \\
\textbf{검증의 정직성} & \textbf{25\%} & 3분할·train-fit·test
한 번·베이스라인·\textbf{에포크 학습 곡선}·validation 튜닝이
보고서에서 추적 가능 \\
수식적 정당화 & 20\% & 배운 공식(파라미터 수, 정규화, ELBO 중
1개 이상)으로 ``왜''를 설명했고, 관측한 숫자와 일치 \\
결과의 정직성과 반성 & 15\% & GBDT보다 못한 부분·학습이 안 된
부분을 숨기지 않고 배운 개념으로 진단 \\
\bottomrule
\end{tabular}
\end{center}

\newpage
\section*{7. 동료평가 (Peer-Review)}
\begin{itemize}
\item \textbf{보고서}: 결론에서 \textbf{``왜 고전
  머신러닝이 아니라 딥러닝인가''}에 구체적 근거로 답해야 한다.
\item \textbf{Peer-Review}: 기말 8주차에 다른 두 팀의 보고서에 대해
  작성 리뷰를 각 1통씩 제출한다. 각 리뷰는 아래 5항목
  체크리스트와, 반박 가능한 질문을 1개 이상 포함해야 한다.
\end{itemize}

\vspace{4pt}
\noindent\textbf{Peer-Review 체크리스트 (5항목 --- 중간 프로젝트의 4항목
+ 1)}
\begin{center}
\small
\begin{tabular}{@{}ll@{}}
\toprule
항목 & 확인할 것 \\
\midrule
문제 정의 & 무엇을 예측/분류하려 했는지 명확한가 \\
방법론 적절성 & 데이터 특성(비정형·불균형·시퀀스)에 맞는
아키텍처를 골랐는가 \\
수식적 정당화 & 요구사항(최소 1개)을 충족하고 실제로 타당한가 \\
결과의 정직성 & 잘 안 된 부분을 숨기지 않고 원인을 분석했는가 \\
\textbf{``왜 딥러닝인가''} & 필수 질문에 구체적 근거로 답했는가 \\
\bottomrule
\end{tabular}
\end{center}

\section*{8. 팀 구성 및 제출 방식}
\begin{itemize}
\item \textbf{팀 구성}: 팀당 인원 및 팀 편성 방식은 추후 공지.
\item \textbf{제출물}: Abstract(기말 2주), 보고서(PDF)와 코드(기말 6주), peer-review 2건(기말 8주).
\item \textbf{제출 방법 및 마감}: 추후 공지 (LMS/과제 제출 시스템
  및 마감 일시는 별도 안내 예정).
\end{itemize}

\end{document}
